% Generated by scripts/blueprint_gen.py; do not edit.
\chapter{The comparison theorem}\label{chap:C}

\begin{definition}[C0 — Resolving data]
\label{C0}
\lean{Cohesive.DgaPresheaf, Cohesive.DgaPresheaf.pre0, Cohesive.PreResolvingData, Cohesive.ResolvingData, Cohesive.PreResolvingData.piecePre, Cohesive.PreResolvingData.pieceSheaf}
\leanok
\uses{A1, A10, A9, S4, S6, S8}
Let X be a topological space. \emph{Pre-resolving data} on X consist of:

\begin{itemize}
\item for every open U a bounded graded-commutative ℂ-dga 𝒜(U) (A1), with restriction maps of dgas for V ⊆ U that are functorial; sections glue uniquely, being of degree q is a local property, and 𝒜\textsuperscript{0} : U ↦ 𝒜\textsuperscript{0}(U) is a sheaf of commutative rings;
\item curvature c = 0 in every 𝒜(U);
\item a sheaf of commutative rings O and ring maps ρ\textsubscript{U} : O(U) → 𝒜\textsuperscript{0}(U), natural in U, with d ∘ ρ\textsubscript{U} = 0;
\item 𝒜\textsuperscript{0} is fine (S1), and every 𝒜\textsuperscript{q} is locally free of finite rank over 𝒜\textsuperscript{0}: each point has an open U and finitely many e\textsubscript{i} ∈ 𝒜\textsuperscript{q}(U) such that, over every V ⊆ U, each β ∈ 𝒜\textsuperscript{q}(V) is uniquely Σ\textsubscript{i} f\textsubscript{i}·e\textsubscript{i}|\textsubscript{V} with f\textsubscript{i} ∈ 𝒜\textsuperscript{0}(V).
\end{itemize}
From these, each 𝒜\textsuperscript{q} is a sheaf of 𝒜\textsuperscript{0}-modules. \emph{Resolving data} add:

\begin{itemize}
\item O → 𝒜\textsuperscript{0} → 𝒜\textsuperscript{1} → ⋯ is a resolution: each ρ\textsubscript{U} is injective, every closed section of 𝒜\textsuperscript{0}(U) lies in the image of ρ\textsubscript{U}, and in degrees q + 1 ≥ 1 every closed form is locally exact;
\item local lifting (LL): for every cohesive module M over 𝒜(X) and x ∈ X there are an open x ∈ U′, a bounded complex W of finitely generated projective O(U′)-modules and an isomorphism in Ho(P\textsubscript{𝒜(U′)}) between M|\textsubscript{U′} and 𝒜(U′) ⊗\textsubscript{O(U′)} W.
\end{itemize}
Deviation: LL is required only for restrictions of global cohesive modules, and K-flatness of 𝒜 over O is not assumed. No compactness or Hausdorff condition is part of the data.
\par\smallskip{\footnotesize Source: blueprint.md §4.}
\end{definition}
\paragraph{Blueprint notes.}
The abstract input, Antweiler's Assumptions 7.2 at h = 0 in the form used here. It is split in two structures: **C0.a \texttt{PreResolvingData}\textbf{ has every field below except the resolution condition and LL; }C0.b \texttt{ResolvingData}** extends it by those two. C1 (the functor J) uses only C0.a, so J can be defined for X without H1 and H2 (see §1).

\begin{itemize}
\item X compact Hausdorff, \texttt{X : Type}; R a sheaf of commutative ℂ-algebras on X.
\item For each open U a bounded graded-commutative dga 𝒜(U) (A1, degrees ≤ N), restriction maps that are dga maps and functorial; each 𝒜\textsuperscript{q} a sheaf of ℂ-modules.
\item 𝒜\textsuperscript{0} a fine sheaf of commutative rings (S1); each 𝒜\textsuperscript{q} a locally free 𝒜\textsuperscript{0}-module of finite rank.
\item (v1.9, audit F1) \textbf{One carrier for A⁰.} The sheaves are not separate data: they are derived from the presheaf of dgas U ↦ (A\textsubscript{U}, 𝒜\textsubscript{U}). 𝒜\textsuperscript{0} is the \texttt{Sheaf \_ CommRingCat} with 𝒜\textsuperscript{0}(U) = \texttt{CommRingCat.of ↥(𝒜\_U 0)} and the restriction maps of the dgas; it is used as a \texttt{RingSheaf} through \texttt{ringOf} (S8). 𝒜\textsuperscript{q}(U) = ↥(𝒜\textsubscript{U} q) as an 𝒜\textsuperscript{0}(U)-module (\texttt{GCDga.piece}). The sheaf conditions are fields. So Γ(X, 𝒜\textsuperscript{0}) and A⁰ = ↥(𝒜\textsubscript{X} 0) are the same type with the same ring structure (prototype P7: \texttt{rfl} at instance transparency), modules over Γ(X, 𝒜\textsuperscript{0}) are \texttt{𝒜 0}-modules (A11 applies to P = Γ(P̃) without transport), and ρ\textsubscript{U} : R(U) →+* ↥(𝒜\textsubscript{U} 0) is the ring map A10 takes.
\item ρ : R → 𝒜\textsuperscript{0} a map of sheaves of rings with d∘ρ = 0 such that 0 → R → 𝒜\textsuperscript{0} → 𝒜\textsuperscript{1} → ⋯ is exact (a resolution).
\item (v1.10) \textbf{Flat curvature: c = 0} in every 𝒜(U) (a field of C0.a). Needed for α(M) to be a complex (𝕏² = c·(−)) and for \texttt{fromComplex} in LL (A10 takes \texttt{hc : c = 0}). It holds for the Dolbeault dga by definition (F10), so F14 fills it; it is not an analytic hypothesis. Curved data (c ≠ 0) is the twisted case Z3 (LATER).
\item \textbf{LL (local lifting).} For every cohesive module M over 𝒜(X) and x ∈ X, there are an open x ∈ U′, a bounded complex W of finitely generated projective R(U′)-modules and a homotopy equivalence in P\textsubscript{𝒜(U′)} between M|U′ (A9) and 𝒜(U′) ⊗\textsubscript{R(U′)} W (A10).
\item \warn{} [A] Assumption 6 is a sheaf version for modules on every open set. C2 and C4 only ever apply LL to restrictions of global modules, so this weaker, sectionwise form suffices; it is also the form of the gauge step in Block's proof of [B, Lemma 4.1.5].
\item \warn{} The first half of [A] Assumption 5 (𝒜 K-flat over R) is \textbf{not} assumed: C5 uses S8 instead. (Its second half, each 𝒜\textsuperscript{q} flat over 𝒜\textsuperscript{0}, follows from local freeness above.)
\end{itemize}

\begin{definition}[C1 — The functor J]
\label{C1}
\lean{Cohesive.PreResolvingData.sheafOf, Cohesive.PreResolvingData.JK, Cohesive.PreResolvingData.J, Cohesive.PreResolvingData.sheafOfOver}
\leanok
\uses{A9, C0, S4}
Let pre-resolving data on X be given (C0), A = 𝒜(X) and R = O. For a cohesive module M = (Q, 𝕏) over A, α(M) is the complex of sheaves of O-modules with α(M)\textsuperscript{n}(V) the degree-n part of the base change M|\textsubscript{V} (A9), i.e. of 𝒜(V) ⊗\textsubscript{𝒜\textasciicircum{}0(V)} ⨁\textsubscript{p} (𝒜\textsuperscript{0}(V) ⊗\textsubscript{A\textasciicircum{}0} Q\textsuperscript{p}) ≅ ⨁\textsubscript{p+q=n} 𝒜\textsuperscript{q}(V) ⊗\textsubscript{A\textasciicircum{}0} Q\textsuperscript{p}, O acting through ρ, with differential 𝕏 of M|\textsubscript{V} (d² = 0 because c = 0).

\begin{itemize}
\item Closed degree-0 maps give chain maps and coboundaries give null-homotopic maps, so M ↦ α(M) defines J\textsubscript{K} : Ho(P\textsubscript{A}) ⥤ K(O) and J : Ho(P\textsubscript{A}) ⥤ D(O), the composite with the localization.
\item For an open U and a cohesive module M over 𝒜(U), α\textsubscript{U}(M) is the same construction on the over-site of U, and α(M)|\textsubscript{U} ≅ α\textsubscript{U}(M|\textsubscript{U}).
\end{itemize}
J lands in the whole derived category D(O); that it lands in D\textsubscript{perf} is C2.
\par\smallskip{\footnotesize Source: blueprint.md §4.}
\end{definition}
\paragraph{Blueprint notes.}
Input: pre-resolving data (C0.a) only; J : Ho(P\textsubscript{A}) ⥤ D(R), into the whole derived category. For a cohesive module M = (Q, 𝔼) over A = 𝒜(X), α(M) is the complex of sheaves of R-modules α(M)\textsuperscript{n}(V) = ⊕\textsubscript{p+q=n} 𝒜\textsuperscript{q}(V) ⊗\textsubscript{𝒜\textasciicircum{}0(X)} Q\textsuperscript{p} with the differential of the base change M|V (A9).

\begin{itemize}
\item \textbf{C1.a} It is a sheaf (Q\textsuperscript{p} is finitely generated projective; S4.c). R acts through ρ, and the differential is R-linear because d∘ρ = 0.
\item \textbf{C1.b} Closed degree-0 maps go to chain maps, homotopies to homotopies, so J : Ho(P\textsubscript{A}) ⥤ K(R) ⥤ D(R). [B, §4.1, the sheaves ℰ\textsubscript{X}]
\item \textbf{C1.c} α(M)|U ≅ α\textsubscript{U}(M|U), naturally (sheaf on U from the cohesive module over 𝒜(U)).
\end{itemize}

\begin{theorem}[C2 — J lands in perfect complexes]
\label{C2}
\lean{Cohesive.ResolvingData.quasiIso_of_gamma_eq, Cohesive.ResolvingData.exists_quasiIso_sheafOfOver_fromComplex, Cohesive.ResolvingData.isPerfect_sheafOf, Cohesive.ResolvingData.perf_J_obj}
\leanok
\uses{C1, C7, S4, S9}
Let resolving data be given on a compact Hausdorff space X.

\begin{itemize}
\item For an open U and a bounded complex W of finitely generated projective O(U)-modules there are a strictly perfect complex L of O|\textsubscript{U}-modules and a quasi-isomorphism L → α\textsubscript{U}(𝒜(U) ⊗\textsubscript{O(U)} W) (A10).
\item For every cohesive module M over 𝒜(X), α(M) is a perfect complex; hence J(M) ∈ D\textsubscript{perf}(X, O).
\end{itemize}
Deviation: X is assumed compact Hausdorff, which the proof uses through the explicit model C7.
\par\smallskip{\footnotesize Source: blueprint.md §4.}
\end{theorem}
\begin{proof}
\leanok
See the Lean proof and the cited items.
\paragraph{Blueprint notes.}
For every cohesive M, α(M) is perfect (S9.b). [A, Lemma 7.2.17] Proof: by LL and C1.c, α(M)|U′ is homotopy equivalent to 𝒜|U′ ⊗\textsubscript{R} W̃, W̃(V) = R(V) ⊗ W strictly perfect; W̃ → 𝒜 ⊗\textsubscript{R} W̃ is a quasi-isomorphism because ρ is a resolution and each W̃\textsuperscript{p} is a direct summand of a finite free module.

\begin{itemize}
\item (v1.11) \textbf{Proof via C7} (the identification above was not itemized). X compact Hausdorff. e : Tot(𝒜(U′) ⊗ W) → Tot(N\textsubscript{W̃}(U′)), a ⊗ (α ⊗ w) ↦ aα ⊗ η(1 ⊗ w̃) with w̃ ∈ Γ(W̃) the section corresponding to w. gr\textsuperscript{0} Ξ\textsubscript{e} is an isomorphism (both sides summands of finite free 𝒜⁰|U′-modules with the same global sections, S4.c on the over-site), so Ξ\textsubscript{e} is a quasi-isomorphism (C7.f). u ≫ Ξ\textsubscript{e} = (W̃ → T(W̃)) (both determined by Γ, S4.c), a quasi-isomorphism by C7.d; so u is one (2-out-of-3).
\item (v1.9, audit F5) \textbf{The roof.} \texttt{IsPerfect} wants α(M)|U′ ← L → P. Take L := W̃, P := W̃, g := id and f := α(h) ∘ u : W̃ → 𝒜|U′ ⊗\textsubscript{R} W̃ → α(M)|U′, where h : 𝒜(U′) ⊗ W → M|U′ is the direction of the LL homotopy equivalence pointing into M|U′ (a homotopy equivalence, so α(h) is a quasi-isomorphism).
\item (v1.9) \textbf{W̃ without fineness.} R|U′ is not fine. W̃ is built on the over-site of U′ from S4.c, whose parts used here (\texttt{exists\_globalSections\_iso}, the Hom bijection) need neither fineness nor compactness; transfer to the over-site through \texttt{TopologicalSpace.Opens.sheafOfModulesEquivOver}, or restate S4.c on any site. The differentials of W lift uniquely by the Hom bijection.
\end{itemize}
\end{proof}

\begin{theorem}[C3 — Hom sheaves of cohesive modules]
\label{C3}
\lean{Cohesive.PreResolvingData.homAlong, Cohesive.PreResolvingData.homTop, Cohesive.PreResolvingData.isSoft_homSh}
\leanok
\uses{A4, S3.b, S4, S7}
Let pre-resolving data be given on X and M₁, M₂ cohesive modules over 𝒜(U₀). For V ⊆ U₀ let ℋ\textsuperscript{n}(M₁, M₂)(V) be the ℂ-linear maps χ : Tot(M₁) → Tot(M₂|\textsubscript{V}) of degree n with χ(a·x) = (−1)\textsuperscript{n|a|} a|\textsubscript{V}·χ(x); by base change these are the degree-n maps M₁|\textsubscript{V} → M₂|\textsubscript{V}.

\begin{itemize}
\item For U₀ = X these form a complex ℋ(M₁, M₂) of sheaves of O-modules with δχ = 𝕏₂ ∘ χ − (−1)\textsuperscript{n} χ ∘ 𝕏₁, and ℋ(M₁, M₂)(X) ≅ Hom\textsuperscript{•\_}\{P\textsubscript{A}\}(M₁, M₂) in each degree.
\item If X is compact Hausdorff, every term ℋ\textsuperscript{n}(M₁, M₂) is soft (S3.b): 𝒜\textsuperscript{0}(X) acts by multiplication. The differential is not 𝒜\textsuperscript{0}-linear.
\end{itemize}
\par\smallskip{\footnotesize Source: blueprint.md §4.}
\end{theorem}
\begin{proof}
\leanok
\uses{Z3.m}
See the Lean proof and the cited items.
\paragraph{Blueprint notes.}
For cohesive M₁, M₂ over A, ℋ(M₁, M₂)(V) = Hom\textsuperscript{•\_}\{P\textsubscript{𝒜(V)}\}(M₁|V, M₂|V). It is a complex of sheaves of R-modules whose terms are 𝒜\textsuperscript{0}-modules (A4.c and S4.c); \warn{} the differential is not 𝒜\textsuperscript{0}-linear (δ(fφ) = df·φ + f·δφ). ℋ(M₁, M₂)(X) = Hom\textsuperscript{•\_}\{P\textsubscript{A}\}(M₁, M₂).

\begin{itemize}
\item \textbf{C3.a} (v1.9, audit F3) \emph{The actions needed for softness.} (i) 𝒜\textsuperscript{0}(X) acts on each term α(M₁)\textsuperscript{i} by R-linear sheaf endomorphisms (multiplication; R-linear because R → 𝒜\textsuperscript{0} is a map of commutative rings), additively, unitally, vanishing on V where f|V = 0; by S7.a′ it acts on each term of 𝓗om\textsubscript{R}(α(M₁), N) for every N. (ii) 𝒜\textsuperscript{0}(X) acts on each term of ℋ(M₁, M₂) (by restriction and multiplication in 𝒜(V)), with the same three properties. Both actions are termwise, not chain maps; these are the inputs of S3.b(i) in C4 step 2.
\end{itemize}
\end{proof}

\begin{theorem}[C4 — Full faithfulness]
\label{C4}
\lean{Cohesive.ResolvingData.bijective_J_map, Cohesive.ResolvingData.quasiIso_Φ, Cohesive.ResolvingData.bijective_J_map_comp, Cohesive.PreResolvingData.Ψ_top, Cohesive.ResolvingData.quasiIso_Ψ_comp, Cohesive.PreResolvingData.exists_isStrictlyGE_ℋ, Cohesive.ResolvingData.bijective_J_map}
\leanok
\uses{C2, C3, S11, S3.b, S7}
Let resolving data be given on a compact Hausdorff space X. Then for all M, N in Ho(P\textsubscript{A}), J : Hom(M, N) → Hom\textsubscript{D(O)}(J M, J N) is bijective. Main steps:

\begin{itemize}
\item Ψ : ℋ(M₁, M₂) → 𝓗om(α(M₁), α(M₂)), χ ↦ α(χ), is a map of complexes, and ℋ(M₁, M₂) is bounded below; at X it sends a closed degree-0 map φ to α(φ).
\item For a quasi-isomorphism r : α(M₂) → I to a bounded-below complex of injectives, r\textsubscript{*} ∘ Ψ : ℋ(M₁, M₂) → 𝓗om(α(M₁), I) is a quasi-isomorphism (locally via LL, then S6).
\item Hence f ↦ J(f) ∘ Q(r) is a bijection Hom\textsubscript{Ho(P\_A)}(M, N) → Hom\textsubscript{D(O)}(α(M), I).
\end{itemize}
\par\smallskip{\footnotesize Source: blueprint.md §4.}
\end{theorem}
\begin{proof}
\leanok
See the Lean proof and the cited items.
\paragraph{Blueprint notes.}
J is fully faithful. [A, Prop 7.2.19]

\begin{enumerate}
\item Choose a bounded-below injective resolution r : α(M₂) → I (S11).
\item Ψ : ℋ(M₁, M₂) → 𝓗om\textsubscript{R}(α(M₁), I), φ ↦ r ∘ α(φ), is a map of complexes of sheaves of R-modules; 𝒜\textsuperscript{0}(X) acts on the terms of both sides (C3.a, S7.a′), hence they are soft (S3.b(i)), but the differentials are not 𝒜\textsuperscript{0}-linear (v1.7). The terms of cone(Ψ) are biproducts of terms of both sides, soft by the closure lemmas of S3.b.
\item Ψ is a quasi-isomorphism of sheaf complexes. Locally (S6), at x take U′ and the homotopy equivalence h : 𝒜(U′) ⊗ W → M₁|U′ from LL. In the square
\end{enumerate}
\begin{verbatim}
   ℋ(M₁,M₂)|U′ ───Ψ───> 𝓗om_R(α(M₁), I)|U′
       │ h^*                  │ α(h)^*
   ℋ(𝒜⊗W, M₂) ───Ψ───> 𝓗om_R(𝒜 ⊗_R W̃, I)
       │ ≅ (A10.a)            │ u^*,  u : W̃ → 𝒜 ⊗_R W̃
   𝓗om_R(W̃, α(M₂)) ──r_*──> 𝓗om_R(W̃, I)
\end{verbatim}
both upper vertical maps are sectionwise homotopy equivalences, u\textsuperscript{\emph{} is a quasi-isomorphism by S7.c (cone(u) is acyclic), and r\textsubscript{}} is one by S7.b. (v1.11) Lean: the square lives on the over-site of U′, so it needs ℋ\textsubscript{U′} (C3's section data with base U′) with ℋ(M₁, M₂)|U′ ≅ ℋ\textsubscript{U′}(M₁|U′, M₂|U′), and Ψ\textsubscript{U′} compatible with Ψ under this iso, \texttt{restrSheafOfIso} and \texttt{homComplexOverIso}; u is C2's map.

\begin{enumerate}
\item Γ(X, Ψ) is a quasi-isomorphism: S3.b(iv) applied to cone(Ψ), bounded below (α(M₁) and ℋ bounded, I bounded below) with soft terms (v1.7; S3 does not apply, see C3).
\item Take H\textsuperscript{0}: Hom\textsubscript{Ho(P\_A)}(M₁, M₂) ≅ H\textsuperscript{0} Hom(α(M₁), I) ≅ Hom\textsubscript{D(R)}(α(M₁), α(M₂)) (S11), and check that the composite is J on morphisms.
\end{enumerate}
\end{proof}

\begin{theorem}[C5 — Essential surjectivity]
\label{C5}
\lean{Cohesive.ResolvingData.exists_J_obj_iso, Cohesive.PreResolvingData.ρSheaf, Cohesive.ResolvingData.exists_flat_model, Cohesive.ResolvingData.isPerfect_sheafBC, Cohesive.ResolvingData.exists_sheafOfMap_quasiIso, Cohesive.ResolvingData.quasiIso_tensorCMap_ρhat, Cohesive.ResolvingData.tensorC_single_unit_iso}
\leanok
\uses{A10, A11, A7, A9, C7, S10, S11, S3, S4, S8, S9}
Let resolving data be given on a compact Hausdorff space X. Every object of D\textsubscript{perf}(X, O) is isomorphic to J(M) for some cohesive module M over 𝒜(X). Main steps:

\begin{itemize}
\item a perfect K receives a quasi-isomorphism from a bounded-above complex K′ of flat O-modules;
\item for such K′ (perfect), the base change 𝒜\textsuperscript{0} ⊗\textsubscript{O} K′ along ρ : O → 𝒜\textsuperscript{0} is perfect;
\item given a strictly perfect complex P̃ of 𝒜\textsuperscript{0}-modules and a quasi-isomorphism P̃ → 𝒜\textsuperscript{0} ⊗\textsubscript{O} K′, there are a cohesive M and a quasi-isomorphism α(M) → 𝒜\textsuperscript{•} ⊗\textsubscript{O} K′ (𝒜\textsuperscript{•} the first factor);
\item for bounded-above, termwise flat K′ (no compactness needed), ρ̂ ⊗ K′ : O[0] ⊗ K′ → 𝒜\textsuperscript{•} ⊗ K′ is a quasi-isomorphism, and O[0] ⊗ K′ ≅ K′.
\end{itemize}
Deviation: K-flat resolutions replace Malgrange's flatness theorem.
\par\smallskip{\footnotesize Source: blueprint.md §4.}
\end{theorem}
\begin{proof}
\leanok
See the Lean proof and the cited items.
\paragraph{Blueprint notes.}
Every object of D\textsubscript{perf}(X, R) is isomorphic to some J(M). [A, Prop 7.2.18], \warn{} K-flat variant.

\begin{itemize}
\item \textbf{C5.0 Tensor input} (the only place where flatness enters). For every perfect K there is a K′ connected to K by quasi-isomorphisms such that (i) K′ → 𝒜 ⊗\textsubscript{R} K′ is a quasi-isomorphism and (ii) 𝒜\textsuperscript{0} ⊗\textsubscript{R} K′ is a globally bounded perfect complex of 𝒜\textsuperscript{0}-modules. Steps 1, 2 and 6 below prove C5.0 with S8; steps 3–5 and 7 use only C5.0. Keeping this interface lets C5 be written before S8 is finished, and keeps open the alternative proof by Malgrange flatness (K′ = K).
\item Let K be perfect. It is globally bounded (S9.d); replace it by a bounded complex (truncations \texttt{CochainComplex.truncGE} ✓, \texttt{truncLE}) and then by a K-flat resolution K′ → K (S8.d).
\item 𝒜\textsuperscript{0} ⊗\textsubscript{R} K′ is a globally bounded perfect complex of 𝒜\textsuperscript{0}-modules (S8.f, S8.g). (v1.9, audit F4) The middle term of each local roof K′|U ← L → P\textsubscript{U} is unbounded; truncate it to τ\textsubscript{≤b}L and resolve by S8.d before applying S8.g, as in S8.g's v1.9 note.
\item By S10 there is a strictly perfect complex P̃ of 𝒜\textsuperscript{0}-modules quasi-isomorphic to it; P̃ is K-projective (S4.a, S11), so the quasi-isomorphism is an honest map f\textsuperscript{0} : P̃ → 𝒜\textsuperscript{0} ⊗\textsubscript{R} K′.
\item Pass to global sections: P = Γ(P̃) is a bounded complex of finitely generated projective 𝒜\textsuperscript{0}(X)-modules (S4.c); Γ(f\textsuperscript{0}) is a quasi-isomorphism (S3); N := the quasi-cohesive module of C5.a. Each 𝒜\textsuperscript{q}(X) is projective, hence flat, over 𝒜\textsuperscript{0}(X) (S4.b, S4.c), so A11 gives a cohesive M on 𝒜(X) ⊗ P and a closed e : M → N extending Γ(f\textsuperscript{0}).
\item Localize e to a map of complexes of sheaves α(M) → 𝒜 ⊗\textsubscript{R} K′ (C5.b); it is a quasi-isomorphism (C5.c).
\end{itemize}
(v1.11) \textbf{C5.a–c are proved via C7} with U = X and K = K′, which replaces the sheaf-tensor arguments below (they relied on a base-change associativity that S8.a does not contain, and on an 𝒜(V)-action on sections of a sheafified tensor product): N := N\textsubscript{K′}(X) (C7.c); A11 with e⁰ = Γ(f⁰) (a quasi-isomorphism by S3: 𝒜⁰ is fine, no boundedness needed); C5.b := Ξ\textsubscript{e} (C7.e); C5.c := C7.f, with gr\textsuperscript{0} Ξ\textsubscript{e} ≅ f⁰ via Γ (S4.c); step 7 goes through T(K′) and C7.d. The text below is kept as the informal picture.

Sub-items (v1.9, audit F2):

\begin{itemize}
\item \textbf{C5.a Quasi-cohesive structure on Γ(𝒜 ⊗\textsubscript{R} K′).} Q\textsuperscript{p} := Γ(𝒜\textsuperscript{0} ⊗\textsubscript{R} K′\textsuperscript{p}) (an 𝒜\textsuperscript{0}(X)-module). This is \emph{not} A10 (K′ is not finitely generated projective, and Γ(𝒜 ⊗\textsubscript{R} K′) ≠ 𝒜(X) ⊗\textsubscript{R(X)} Γ(K′)). Instead: 𝒜\textsuperscript{q} ⊗\textsubscript{R} K′\textsuperscript{p} ≅ 𝒜\textsuperscript{q} ⊗\textsubscript{𝒜\textasciicircum{}0} (𝒜\textsuperscript{0} ⊗\textsubscript{R} K′\textsuperscript{p}) (base-change associativity, S8.a), and Γ of the right side is 𝒜\textsuperscript{q}(X) ⊗\textsubscript{𝒜\textasciicircum{}0(X)} Q\textsuperscript{p} by S4.d (𝒜\textsuperscript{q} is a summand of a finite free 𝒜\textsuperscript{0}-module, F11 and S4.b). 𝕏 is Γ of the sheaf operator 𝕏(b ⊗ k) = db ⊗ k + (−1)\textsuperscript{|b|} b ⊗ d\textsubscript{K′}k (well defined because d∘ρ = 0), transported along these isomorphisms; Leibniz and 𝕏² = 0 hold sectionwise, hence after transport. (v1.10, sign fix) The Koszul sign is essential: without it Leibniz fails whenever 𝒜 ≠ 𝒜\textsuperscript{0} (𝕏(a·(b⊗k)) = da·(b⊗k) + (−1)\textsuperscript{|a|} a·𝕏(b⊗k) needs (−1)\textsuperscript{|a|+|b|} on ab ⊗ dk). This operator is the differential of \texttt{tensorC R 𝒜\textasciicircum{}• K′} with 𝒜\textsuperscript{•} as the \textbf{first} factor (Mathlib puts the sign on the second factor), as fixed in conventions §2. Its (0,0)-component is Γ(𝒜\textsuperscript{0} ⊗\textsubscript{R} d\textsubscript{K′}), so A11's target complex (Q, 𝕏\textsuperscript{0}) is Γ(𝒜\textsuperscript{0} ⊗\textsubscript{R} K′).
\item \textbf{C5.b The sheaf map.} On V, e\textsubscript{V} : 𝒜(V) ⊗\textsubscript{𝒜\textasciicircum{}0(X)} P → Γ(V, 𝒜 ⊗\textsubscript{R} K′), a ⊗ p ↦ a·res\textsubscript{V}(e(1 ⊗ p)). (v1.10) No sign: e is closed of degree 0, and a degree-0 map satisfies φ(a·m) = a·φ(m) (conventions §2, Hom\textsuperscript{n} with n = 0). For a map of degree n the generator extension carries (−1)\textsuperscript{n|a|}. N|V is \emph{not} the base change of N, so A9 does not supply the map; α(M)(V) is A9's base change of M, and e\textsubscript{V} is defined on generators. It is a chain map because res\textsubscript{V} is a dga map and e is closed; it is compatible with restriction, so it defines α(M) → 𝒜 ⊗\textsubscript{R} K′ (S4.c Hom property on each α(M)\textsuperscript{n}, which is a finite sum of summands of finite free 𝒜\textsuperscript{0}-modules).
\item \textbf{C5.c Quasi-isomorphism.} Form-degree filtration on both sides as complexes of sheaves of abelian groups; A7 applies as stated (\texttt{quasiIso\_of\_finite\_filtration} is proved in any abelian category). On gr\textsuperscript{b} the map is 𝒜\textsuperscript{b} ⊗\textsubscript{𝒜\textasciicircum{}0} f\textsuperscript{0} (up to the shift sign of A12), a quasi-isomorphism because 𝒜\textsuperscript{b} is a summand of a finite free 𝒜\textsuperscript{0}-module (tensoring with it preserves quasi-isomorphisms; flat by S8.b, or directly).
\item K′ → 𝒜 ⊗\textsubscript{R} K′ is a quasi-isomorphism (𝒜\textsuperscript{•} first, as in C5.a): its cone is cone(ρ) ⊗\textsubscript{R} K′, acyclic because ρ is a resolution and K′ is K-flat.
\item Hence α(M) ≅ 𝒜 ⊗\textsubscript{R} K′ ≅ K′ ≅ K in D(R).
\end{itemize}
\end{proof}

\begin{theorem}[C6 — Comparison theorem for resolving data]
\label{C6}
\lean{Cohesive.ResolvingData.comparison}
\leanok
\uses{C2, C4, C5}
Let resolving data be given on a compact Hausdorff space X, with A = 𝒜(X). Then the functor J : Ho(P\textsubscript{A}) ⥤ D(O) is full and faithful, and its essential image is exactly D\textsubscript{perf}(X, O).

This is the abstract comparison theorem, for any resolving data (C0). Its application to the Dolbeault data is F15 (\texttt{blockPerf}), which is the proof of M.
\par\smallskip{\footnotesize Source: blueprint.md §4.}
\end{theorem}
\begin{proof}
\leanok
See the Lean proof and the cited items.
\paragraph{Blueprint notes.}
For resolving data, J : Ho(P\textsubscript{𝒜(X)}) ⥤ D(R) is full and faithful (C4) with essential image exactly D\textsubscript{perf}(X, R) (C2, C5); equivalently its corestriction to D\textsubscript{perf} is an equivalence. [A, Thm 7.2.20 at h = 0], [CHL, Thm 7.10] (for twisted modules) and [CHL, Cor 7.13] (the same for cohesive modules, via [CHL, Cor 3.4]) (full faithfulness: [CHL, Prop 7.4]).
\end{proof}

\begin{theorem}[C7 — The explicit model T(K) of 𝒜 ⊗\textsubscript{R} K]
\label{C7}
\lean{Cohesive.PreResolvingData.exists_TC_iso_tensorC, Cohesive.PreResolvingData.quasiIso_Ξ, Cohesive.PreResolvingData.Conn, Cohesive.PreResolvingData.Conn.N, Cohesive.PreResolvingData.TSh, Cohesive.PreResolvingData.Conn.TC, Cohesive.PreResolvingData.Ξ}
\leanok
\uses{A3, A6, A7, A9, S4, S8}
Let pre-resolving data be given on a compact space X, U ⊆ X open, and F a complex of 𝒜\textsuperscript{0}|\textsubscript{U}-modules on the over-site of U. A \emph{flat connection} on F is, for every V, generator data θ\textsubscript{V} on the 𝒜\textsuperscript{0}(V)-modules F\textsuperscript{p}(V) of total degree 1, flat (𝕏\textsubscript{V²} = c = 0), with form-degree-0 part d\textsubscript{F}, and compatible with restriction. Given one:

\begin{itemize}
\item N(V) = (⨁\textsubscript{p} F\textsuperscript{p}(V), 𝕏\textsubscript{V}) is a quasi-cohesive module over 𝒜(V);
\item T(F)\textsuperscript{n}(V) = (𝒜(V) ⊗\textsubscript{𝒜\textasciicircum{}0(V)} ⨁\textsubscript{p} F\textsuperscript{p}(V))\textsuperscript{n} is a sheaf of O|\textsubscript{U}-modules, and with 𝕏 a complex T;
\item a closed degree-0 map e : M → N(U) from a cohesive module over 𝒜(U) induces a map of complexes Ξ\textsubscript{e} : α\textsubscript{U}(M) → T.
\item For K a complex of O|\textsubscript{U}-modules and F = 𝒜\textsuperscript{0} ⊗\textsubscript{O} K, with a connection satisfying ∇(1 ⊗ k) = 0 (which exists), T ≅ 𝒜\textsuperscript{•}|\textsubscript{U} ⊗ K, and under this and K ≅ O|\textsubscript{U}[0] ⊗ K the map K → T, k ↦ 1 ⊗ (1 ⊗ k), is ρ̂|\textsubscript{U} ⊗ K.
\end{itemize}
\par\smallskip{\footnotesize Source: blueprint.md §4.}
\end{theorem}
\begin{proof}
\leanok
\uses{A11, C1}
See the Lean proof and the cited items.
\paragraph{Blueprint notes.}
Used by C2, C4 step 3 and C5.a–c in place of the sheafified tensor product, whose sections are not computable. Setting: resolving data, X compact Hausdorff, an open U, the over-site of U. For a complex K of R|U-modules let F := 𝒜⁰ ⊗\textsubscript{R} K (S8 \texttt{sheafBC} along ρ), a complex of 𝒜⁰|U-modules, and T(K)\textsuperscript{n}(V) := degree-n part of 𝒜(V) ⊗\textsubscript{𝒜⁰(V)} ⊕\textsubscript{p} F\textsuperscript{p}(V) = ⊕\textsubscript{q} 𝒜\textsuperscript{q}(V) ⊗\textsubscript{𝒜⁰(V)} F\textsuperscript{n−q}(V), with plain tensor products of sections (A3's \texttt{Tot} over 𝒜(V), generators F\textsuperscript{p}(V)).

\begin{itemize}
\item \textbf{C7.a Sheaf.} T(K)\textsuperscript{n} is a sheaf of R|U-modules: 𝒜\textsuperscript{q}|U is a summand of a finite free 𝒜⁰|U-module (S4.b on X, restricted), so T(K)\textsuperscript{n} is sectionwise a retract of a finite power of ⊕\textsubscript{q} F\textsuperscript{n−q}, compatibly with restriction.
\item \textbf{C7.b Connection.} ∇ : F\textsuperscript{p} → 𝒜¹ ⊗\textsubscript{𝒜⁰} F\textsuperscript{p} (target a sheaf by C7.a) is the sheafification of α ⊗ k ↦ dα ⊗ η(1 ⊗ k) on the presheaf tensor 𝒜⁰(V) ⊗\textsubscript{R(V)} K\textsuperscript{p}(V), η the unit of sheafification (R-linear, well defined since d∘ρ = 0). ∇(αf) = dα ⊗ f + α∇f, ∇ commutes with d\textsubscript{F}, and ∇² = 0 (in 𝒜² ⊗ F): checked where f comes from the presheaf, then everywhere (η is locally surjective; the target is separated).
\item \textbf{C7.c Differential.} 𝕏\textsubscript{V}(a ⊗ f) = da ⊗ f + (−1)\textsuperscript{|a|}(a·∇f + a ⊗ d\textsubscript{F} f) (Koszul sign of v1.10; ∇ extended to 𝒜 ⊗ F by Leibniz). Leibniz holds by construction, 𝕏² = 0 from C7.b and c = 0, and 𝕏 commutes with restriction. Hence (F\textsuperscript{•}(V), 𝕏\textsubscript{V}) is a quasi-cohesive module N\textsubscript{K}(V) over 𝒜(V) (A3; generators not finitely generated), the restrictions are closed degree-0 maps semilinear along res, and T(K) is a complex of sheaves of R|U-modules.
\item \textbf{C7.d Comparison with tensorC.} The presheaf map 𝒜\textsuperscript{q}(V) ⊗\textsubscript{R(V)} K\textsuperscript{p}(V) → T(K)(V), a ⊗ k ↦ a ⊗ η(1 ⊗ k), is locally bijective (a retract of a finite power of η), so it induces an isomorphism tensorC(𝒜\textsuperscript{•}|U, K) ≅ T(K) of complexes (differentials compared on a ⊗ k). With C5 step 6 on any site: for K bounded above with flat terms, the map K → T(K), k ↦ 1 ⊗ η(1 ⊗ k), is a quasi-isomorphism.
\item \textbf{C7.e Comparison map.} For M cohesive over 𝒜(U) and e : Tot(M) → Tot(N\textsubscript{K}(U)) closed of degree 0, Ξ\textsubscript{e} : α\textsubscript{U}(M) → T(K) is on V the A9 extension (\texttt{extAlong}) of res\textsubscript{V} ∘ e; it is a map of complexes of sheaves (closedness, C7.c).
\item \textbf{C7.f Quasi-isomorphism criterion.} Filter α\textsubscript{U}(M) and T(K) by form degree (A6 \texttt{filTot}, sectionwise; subsheaves by C7.a and \texttt{degree\_local}), with short exact sequences sectionwise (A12 \texttt{filSC\_shortExact}) and A7 in sheaves. gr\textsuperscript{b} Ξ\textsubscript{e} is a retract of a finite power of gr\textsuperscript{0} Ξ\textsubscript{e} (the splitting of 𝒜\textsuperscript{b} over 𝒜⁰), so Ξ\textsubscript{e} is a quasi-isomorphism as soon as gr\textsuperscript{0} Ξ\textsubscript{e} : (V ↦ 𝒜⁰(V) ⊗\textsubscript{𝒜⁰(U)} Q\textsuperscript{•}) → F is one.
\end{itemize}
\end{proof}
